% Part: turing-machines
% Chapter: machines-computations
% Section: turing-machines

\documentclass[../../../include/open-logic-section]{subfiles}

\begin{document}

\olfileid{tur}{mac}{tur}
\olsection{Mesin Turing}

\begin{explain}
Definisi formal mesin Turing katon abstrak,
nanging sejatine prasaja: definisi iki
mung nyawijikake kabeh informasi kang
dibutuhake kanggo nemtokake cara makaryane
mesin Turing dadi siji struktur matematis.
Informasi iku kalebu (1) kahanan apa wae
kang bisa dialami mesin, (2) simbol apa
wae kang diidinake ana ing pita, (3)
ing kahanan apa mesin kudu diwiwiti,
lan (4) apa wae instruksi mesin.
\end{explain}

\begin{defn}[Mesin Turing]
\emph{Mesin Turing} $M$ yaiku tupel
$\langle Q, \Sigma, q_0,
\delta\rangle$ kang dumadi saka
\begin{enumerate}
\item himpunan \emph{kahanan}~$Q$ kang winates,
\item \emph{alfabet} $\Sigma$ winates kang ngemot
  $\TMendtape$ lan $\TMblank$,
\item \emph{kahanan wiwitan}~$q_0 \in Q$,
\item \emph{himpunan instruksi}~$\delta\colon Q \times \Sigma
  \pto Q \times \Sigma \times \{\TMleft, \TMright, \TMstay\}$ kang winates.
\end{enumerate}
Fungsi parsial~$\delta$ uga diarani
\emph{fungsi transisi} saka~$M$.
\end{defn}

\begin{explain}
Kita nganggep pita tanpa wates mung
ing siji arah. Mula migunani yen simbol
mligi~$\TMendtape$ ditetepake minangka
panandha pucuk kiwa pita. Iki nggampangake
program mesin Turing mangerteni nalika
program mau ``meh'' metu saka pita.
Kita bisa nganggep simbol iki ora tau
ditimpa, yaiku $\delta(q,\TMendtape) = \tuple{q', \TMendtape, x}$
yen $\delta(q,\TMendtape)$ ditegesi.
Sawetara buku ajar nganggo anggepan iki;
kita ora. Nalika mbangun mesin Turing,
kowe mung perlu ngati-ati supaya mesin
ora nate nimpa~$\TMendtape$. Kajaba iku,
ana kahanan nalika idin kanggo nimpa
simbol iki menehi keluwesan kang migunani.
\end{explain}

\begin{ex}
\emph{Mesin Genep:} Sacara formal,
mesin genep iku rangkap papat
$\tuple{Q, \Sigma, q_0, \delta}$ kanthi
\begin{align*}
Q & = \{ q_0, q_1 \} \\
\Sigma & = \{ \TMendtape, \TMblank, \TMstroke \}, \\
\delta(q_0, \TMstroke) & = \tuple{q_1, \TMstroke, \TMright},\\
\delta(q_1, \TMstroke) & = \tuple{q_0, \TMstroke, \TMright},\\
\delta(q_1, \TMblank)  & = \tuple{q_1, \TMblank, \TMright}.
\end{align*}
\end{ex}

\end{document}
